\subsection{Uniqueness Proof for Homogeneous Means} \label{sect-HLP_proof}

In this section we will prove the Uniqueness Theorem for Homogeneous means, which states that the power means \footnote{In some literature this is also referred to as the Hölder mean or the generalized means} (including the geometric mean as a limit case) are the only homogeneous quasi-arithmetic means.
Before beginning the proof, we will first define these means.

\begin{defi}[Quasi Arithmetic Means and Homogeneity]
  Let $\phi : [0,\infty) \rightarrow [0,\infty)$ be continuous and strictly increasing, $\ve{a} = (a_1, \dots, a_n)$ and $\ve{w} = (w_1, \dots, w_n)$ with $a_i, w_i \geq 0$ for $i=1, \dots, n$.
  Then the quasi-arithmetic $\phi$-mean of $\ve{a}$ with weights $\ve{w}$
  $$
  \mean{\ve{a}} = \phi^{-1} \left ( \sum_{i=1}^{n} w_i \phi(a_i) \right)
  $$

  The mean is said to be \emph{homogeneous}, if
  $$
  \mean{t \cdot \ve{a}} = t \cdot \mean{\ve{a}}
  $$
  holds for all $t > 0$.
\end{defi}

\begin{defi}[Power mean]
  For $n \in \mathbb{N}, r > 0$ and $\ve{a} = (a_1, \dots,  a_n)$ with $a_i > 0$ for all $i=1, ..., n$ we define the power means $\meanr{\ve{a}}$ as
  $$ 
  \meanr{\ve{a}} = \left ( \frac{1}{n} \sum_{i=1}^{n} a_i^r \right)^{1/r}.
  $$

  For $r=0$ we can define $\meanr[0]{\ve{a}}$ as the geometric mean
  $$
  \meanr[0]{\ve{a}} = \sqrt[n]{\prod_{i=1}^n a_i}
  $$
  since it is obtained as a limit case for $r \rightarrow 0$ \parencite[p. 176]{handbookMeans}.
\end{defi}

\begin{rem}[Power Means are Quasi-Arithmetic]
By setting $\phi(x) = x^r$ for $r>0$ or $\phi(x) = \log(x)$ we obtain the power means for $r>0$ and $r=0$ respectively.
\end{rem}


\begin{thm}[Characterization of Equal Means] \label{thm-quasi-mean-equal-cond}
  In order that
  $$
  \mean{\ve{a}}  = \mean[\Psi]{\ve{a}}
  $$
  for all $\ve{a}$ and $\ve{w}$, it is necessary and sufficient that 
  $$
    \phi = \alpha \cdot  \psi + \beta
  $$
  where $\alpha$ and $\beta$ are constants and $\alpha \neq 0$.
\end{thm}


\begin{thm}[$\meanrempty$ are the only homogeneous means $\meanempty$] \label{thm-power-mean-uniqueness}
      Suppose that $\phi$ is continuous and strictly increasing in the open interval $(0,\infty)$ and that the mean defined by $\phi$ is homogeneous, i.e.
      $$ 
      \mean{t \ve{a} } = t \mean{\ve{a}}
      $$
      for all positive $\ve{a}, \ve{w}$ and $t$. Then $\mean{ \ve{a}}$ is $\meanr{\ve{a} }$.

\end{thm}


